\scanpage{46}
\mspage{2}
\[
(8)\qquad g_E=\underbrace{g_{V_J(\mathbb R)}}_{\text{tautol.}}\oplus\operatorname{sg}(g)\operatorname{id}_\Delta
\]
L'opération de $\tau$ est donnée par
\[
(9)\qquad\tau_E=\operatorname{id}_{V_J(\mathbb R)}\oplus(-\operatorname{id}_\Delta)
\]
Je pose
\[
(10)\qquad\Sigma^*=\Sigma\setminus J=\Sigma-\{R_{\deleted{\unsure{j}}\inserted{i}}\mid i\in I\}
\]
et je m'intéresse au groupoïde fondamental de $\Sigma^*$, comme exercice préliminaire de syntaxe galoisienne, avant de regarder des situations \inserted{un peu} plus complexes (le triangle équilatéral remplacé par d'autres configurations riches en symétries sur la sphère, voire sur d'autres surfaces que la sphère, ou celle-ci étant remplacée par le fibré de Hopf, en vue des applications au calcul des groupoïdes de Teichmüller...).

Il est commode d'introduire une compactification de $\Sigma^*$, pour disposer ainsi de ``pts-base'' ``proches'' \unsure{desdits} $R_i$.
\margin{45}
